\documentclass[11pt]{article}
\usepackage[margin=0.85in]{geometry}
\usepackage{lmodern}
\usepackage{microtype}
\usepackage{enumitem}
\usepackage{listings}
\usepackage[T1]{fontenc}

\lstset{
  basicstyle=\ttfamily\small,
  columns=fullflexible,
  keepspaces=true,
  showstringspaces=false,
  breaklines=true,
  frame=single,
  xleftmargin=0.4em,
  xrightmargin=0.4em,
  aboveskip=0.5em,
  belowskip=0.5em
}

\pagestyle{empty}
\setlength{\parskip}{0.32em}
\setlength{\parindent}{0pt}
\setlist[enumerate,1]{leftmargin=1.4em, itemsep=0.65em, topsep=0.25em}
\setlist[enumerate,2]{leftmargin=1.5em, itemsep=0.28em, topsep=0.2em, label=\alph*.}
\setlist[itemize]{leftmargin=1.3em, itemsep=0.12em, topsep=0.15em}

\begin{document}

\begin{center}
{\LARGE\bfseries Final exam --- Sample 1}\\[0.35em]
DATA 412/612
\end{center}

\vspace{0.2em}

Practice exam, questions only. Weeks 1--14. The newest material
is lists and \texttt{purrr} (\texttt{map}, \texttt{map2},
\texttt{keep}, and split-plus-map). Earlier weeks stay in scope.

Put your answers in an R Markdown document and keep the code visible
(\texttt{knitr::opts\_chunk\$set(echo = TRUE)} in the setup chunk).
You may use your notes and the course website.
You may not use Google, any other website, or an LLM.
Do not discuss the exam with anyone else.
Load the tidyverse.

\begin{enumerate}
\item \textbf{Lists} (6 pts).
Build a named list \texttt{fit} with components \texttt{n = 10} and
\texttt{center = c(1, 2)}.
Extract \texttt{center} so the result is the numeric vector, not a
length-one list.
Then replace \texttt{n} with \texttt{12} without rebuilding the list.

\item \textbf{\texttt{map}} (8 pts).
Keep only the numeric columns of \texttt{mpg} and, with
\texttt{map\_dbl()}, return the mean of each.
Pass \texttt{na.rm = TRUE}.
The function argument must be a function, not a call.

\item \textbf{\texttt{map2}} (8 pts).
\texttt{mu} is \texttt{c(0, 10, 20)} and \texttt{sigma} is
\texttt{c(1, 1, 2)}.
Use \texttt{map2()} and \texttt{rnorm} to draw 5 values from each
pair. Set the seed to \texttt{612} first.
Return a list of three numeric vectors.

\item \textbf{Split, then map} (10 pts).
Split \texttt{mpg} by \texttt{drv}.
On each piece, fit \texttt{hwy} on \texttt{displ} and return the
slope on \texttt{displ}.
The result should be a named numeric vector, one slope per drive.

\item \textbf{One pipeline from earlier weeks} (8 pts).
Load \texttt{nycflights13}.
Add the airline name onto \texttt{flights}.
For each name, report the number of flights and the mean arrival
delay, ignoring missing delays.
Keep names with at least 1000 flights.
Make a bar plot of that mean delay, with the names ordered by the
mean. Title the plot.
\end{enumerate}

\end{document}
